多模态大语言模型（MLLM）仍易受幻觉影响：文本流畅但与视觉证据脱节、虚构物体、以及语言先验驱动的错误。本文论证：当学习模块消费宿主模型原生隐状态，或直接在宿主词表几何中写入修正时，严格的“可热插拔”学习控制是不可能的。我们形式化这一\emph{通用性边界}，并提出\textbf{UniGround}：一种在附着阶段无需再训练的解码时外挂（sidecar），其唯一学习核心$\Psi_{\mathrm{univ}}$运行于由冻结的公开视觉/文本编码器、解码得到的候选片段与确定性附着逻辑所构成的共享外部语义空间。UniGround通过\emph{候选驱动的视觉共振}使每个候选片段唤醒局部区域证据，并以带安全门控的保守 logit 重分配仅扰动宿主的活跃 Top-$M$ 解码前沿。我们给出一套评测协议，区分模块级验证、跨模型热插拔有效性、全局先验基线、结构保持与显式失效边界。
